% !PRINTED 496
\weberpage{496}
%% [page furniture, not body text -- source-confirmed 2026-08-10]
%% Nineteenth Chapter. \hfill \S\ 133.

Since 1, 2, 4 are the quadratic residues of 7, and 1, 3, 4, 5, 9 the
quadratic residues of 11, these two sums can be determined according to
Bd.~I, \S\ 179, and in the two cases they give:
\[
i\sqrt{7} \quad \text{and} \quad -i\sqrt{11},
\]
so that
\weq{18}{\begin{aligned}
w_\nu &= q^{-\frac{1}{7}} \varrho^{3\nu} + \cdots, & n &= 7, \\
w_\nu &= -q^{-\frac{5}{22}} \varrho^{5\nu} + \cdots, & n &= 11.
\end{aligned}}
From this an upper limit can be derived, up to which the powers of
$f(\omega)$ in $A_\nu$ can at most rise, provided one observes further
that the $A_\nu$ are representable as integral rational (and
symmetric) functions of the $w_\nu$.

For $n=7$ it follows in this way, since all exponents of $f(\omega)$
must be divisible by 24, that $A_1$, $A_2$, $A_3$, $A_4$, $A_5$, $A_6$
are constants, while $A_7$ contains only the first power of
$f(\omega)^{24}$, and that with the coefficient $-1$ [by (18)].

For $n=11$, if $f(\omega)^{12\lambda}$ is the highest power of
$f(\omega)$ occurring in $A_\nu$, then
\[
\lambda \leqq \frac{5\nu}{11}
\]
must hold, and moreover $\lambda$ must be even or odd at the same time
as $\nu$. Hence $A_2$, $A_4$ are constant, $A_6$, $A_8$ contain the
highest power $f(\omega)^{24}$, and $A_{10}$ contains $f(\omega)^{48}$
as well. Since the $A$ with odd index contain only odd powers of
$f(\omega)^{12}$, we have $A_1 = 0$, while $A_3$, $A_5$ contain
$f(\omega)^{12}$, $A_7$, $A_9$ contain up to $f(\omega)^{36}$, and in
$A_{11}$ the function $f(\omega)$ rises to the 60th power.

For the case $n=7$ we will now determine the constants completely by
means of complex multiplication, and for this the consideration of the
values of the $w_\nu$ for $\omega = i$ suffices.

For $\omega = i$ we get
\[
u = f(i) = \sqrt[4]{2},
\]
and if we introduce this value into equation (2) and put
\[
\sqrt[4]{2}\, v = x
\]
we obtain:
\weq{19}{x^8 - 4x^7 + 28x^4 - 32x + 16 = 0.}
The left-hand side of this equation is the square of the expression:
\weq{20}{x^4 - 2x^3 - 2x^2 - 4x + 4,}
%CONT>>
